\documentclass[a4paper, 12pt]{article}

\usepackage[english, russian]{babel}
\usepackage[T2A]{fontenc}
\usepackage[utf8]{inputenc}
\usepackage{amsmath, amsfonts, amssymb, amsthm, mathtools}
\usepackage{indentfirst}
\usepackage{geometry}

\geometry{top=25mm}
\geometry{bottom=30mm}
\geometry{left=20mm}
\geometry{right=20mm}
\linespread{1}
\setlength{\parindent}{20pt}
\setlength{\parskip}{12pt}

\begin{document}
\section{Основы физики}
    \subsection{Аксиомы}
        В рамках данной работы будет использоваться максимально аксиоматический подход, 
        поэтому в первой главе займемся определением базовых аксиом и выводом следствий из них

        Мы будем рассматримавать двумерную евклидову плоскость (далее иногда будем называть её пространством). Прежде 
        всего положим, что она 
        \textbf{однородна} "--- во всех точках пространства действуют одинаковые законы физики и 
        перемещение системы целиком из одной области пространства в другую не меняет её поведения. 
        Также положим, что наша плоскость \textbf{изотропна} "--- физические свойства не зависят от направления
        и при повороте системы на любой угол её свойства не изменятся. Время течет одинаково во всех системах отсчета.

        \textbf{Материальная точка} "--- простейшее тело, которому соответствует положительная скалярная величина ($m>0$),
        называемая \textbf{массой} и положение в пространстве.
\end{document}